% Figure: the pole-zero map of a rational transfer function read as a
% system preview. Complex s-plane with the imaginary axis as the
% stability boundary; a conjugate pole pair in the left half-plane
% (decaying oscillation), a real pole on the negative axis (decaying
% mode), and a zero. The left half-plane is shaded as the stable region.
\begin{figure}[htbp]
\centering
\resizebox{0.62\linewidth}{!}{%
\begin{tikzpicture}[scale=1.0,
  ax/.style={->, draw=engblack, thick},
  stab/.style={fill=engblue!8},
  pole/.style={engred, very thick},
  zero/.style={draw=engblue, very thick},
  ann/.style={font=\footnotesize, align=center},
]
  % Stable (left half-plane) shading.
  \fill[stab] (-4.2, -3.0) rectangle (0, 3.0);
  % Axes: real (sigma) horizontal, imaginary (j omega) vertical.
  \draw[ax] (-4.4, 0) -- (2.2, 0) node[below right, ann] {$\operatorname{Re}(s)=\sigma$};
  \draw[ax] (0, -3.2) -- (0, 3.3) node[left, ann] {$\operatorname{Im}(s)=\mathrm{j}\omega$};
  % Stability-boundary label on the imaginary axis.
  \node[ann, anchor=west, text=engblack] at (0.15, -2.7)
    {stability\\boundary};
  \node[ann, text=engblue!70!black] at (-2.7, 2.5) {stable region\\(decaying modes)};
  \node[ann, text=engred!80!black] at (1.2, 2.5) {unstable\\(growing modes)};
  % Conjugate pole pair: sigma = -1.2, omega = +/- 2.0.
  \draw[pole] (-1.5, 1.7) -- (-0.9, 2.3);
  \draw[pole] (-1.5, 2.3) -- (-0.9, 1.7);
  \draw[pole] (-1.5, -1.7) -- (-0.9, -2.3);
  \draw[pole] (-1.5, -2.3) -- (-0.9, -1.7);
  \node[ann, anchor=west, text=engred!80!black] at (-0.75, 2.0)
    {pole pair\\$\sigma\pm\mathrm{j}\omega_{d}$};
  % Real pole on the negative real axis.
  \draw[pole] (-3.3, -0.3) -- (-2.7, 0.3);
  \draw[pole] (-3.3, 0.3) -- (-2.7, -0.3);
  \node[ann, below, text=engred!80!black] at (-3.0, -0.4) {real pole};
  % Zero on the negative real axis (open circle).
  \draw[zero] (-0.55, 0) circle (0.3);
  \node[ann, above, text=engblue!70!black] at (-0.55, 0.4) {zero};
  % Distance annotations: real part = decay rate, imag part = ring freq.
  \draw[engblack!50, thin, dashed] (-1.2, 0) -- (-1.2, 2.0);
  \draw[engblack!50, thin, dashed] (0, 2.0) -- (-1.2, 2.0);
  \node[ann, anchor=north, text=engblack!70] at (-1.2, -0.12)
    {$\sigma$ sets decay};
  \node[ann, anchor=east, text=engblack!70] at (-0.06, 2.0)
    {$\omega_{d}$ sets ring};
\end{tikzpicture}%
}%
\caption[The pole-zero map read as a system preview]{The poles
(crosses) and zero (circle) of a rational transfer function plotted in
the complex $s$-plane. The imaginary axis is the stability boundary:
poles in the shaded left half-plane are decaying modes, poles in the
right half-plane are growing modes, poles on the axis are sustained
oscillations. For the conjugate pole pair the real part $\sigma$ sets
the decay rate of the mode and the imaginary part $\omega_{d}$ sets the
frequency at which it rings. Reading stability and dynamics off the
location of the poles is the first move of control engineering, and it
is available to the engineer who can find the roots of a rational
function's denominator.}
\label{fig:vol02:ch02:pole-zero-map}
\end{figure}
